% test-005.tex -- comando \spmoney: tabla completa de monedas y alias
%
% Compilar:  lualatex-dev test-005.tex
% Validar:   show-pdf-tags --xml test-005.pdf | rnv-wrapp
%            verapdf --flavour ua2 --format text test-005.pdf
%
% Generado a partir de la tabla de monedas del manual de usuario: una
% entrada por cada simbolo, alias y codigo ISO. Cada caso es un parrafo
% numerado, para poder referirse a el al escuchar el PDF. Los comentarios
% "doc:" indican la lectura que documenta la tabla. Se usa siempre la
% cifra 500.
%
% La fila del yuan repite el simbolo del yen en la columna Entrada, que es
% la misma entrada: ese caso se omite y el yuan se prueba por sus alias.
%
% Los simbolos del won, el sequel, la rupia, el rublo, la lira y la grivna
% no estan en Latin Modern Math (la fuente matematica por defecto), por
% eso esos casos usan money-math=false, que imprime el simbolo con la
% fuente de texto. La fuente de texto es Libertinus Serif, con el respaldo
% de glifos que usa el manual (bloque luacode copiado de spintent.dtx).
% Los codigos ISO en mayusculas se imprimen con \mathrm y no la necesitan.
\DocumentMetadata
  {
    lang=es-CL,
    pdfversion=2.0,
    pdfstandard=ua-2,
    tagging=on,
    tagging-setup={math/setup={mathml-SE}},
  }
\documentclass{article}
\usepackage[spanish]{babel}
% fuente de texto: Libertinus Serif con el mismo respaldo de glifos del
% manual (copiado tal cual de spintent.dtx)
\usepackage{luacode}
% Para documentación y ejemplos
\begin{luacode*}
local SCALE = 0.85
local fallback_cache = {}
local lookup_cache = {}

luatexbase.add_to_callback("luaotfload.patch_font", function(fontdata, specification, id)
  if type(fontdata) ~= "table" then return end
  local meta = fontdata.shared and fontdata.shared.rawdata and fontdata.shared.rawdata.metadata
  local familyname = meta and meta.familyname
  if type(familyname) ~= "string" then return end
  if familyname:lower():find("libertinus serif", 1, true) then
    local size = fontdata.size
    local key  = tostring(size)
    local fb_id = fallback_cache[key]
    if not fb_id then
      local path = kpse.find_file("Carlito-Regular.ttf", "truetype fonts")
      if path then
        local scaled_size = math.floor(size * SCALE)
        local ok, result = pcall(luaotfload.define_font, "[" .. path .. "]", scaled_size)
        if ok then
          fb_id = result
          fallback_cache[key] = fb_id
        end
      end
    end
    if fb_id then
      local lookup = lookup_cache[fb_id]
      if not lookup then
        local fb_font = font.getfont(fb_id)
        lookup = {}
        for uni, _ in next, fb_font.characters do
          lookup[uni] = fb_id
        end
        lookup_cache[fb_id] = lookup
      end
      fontdata.fallback_lookup = lookup
    end
  end
end, "shekel_fallback_patch")
\end{luacode*}
\usepackage[nomath,mono=false]{libertinus-otf}
\usepackage{unicode-math}
\usepackage{spintent}
\usepackage{hyperref}
\hypersetup
  {
    pdfdisplaydoctitle = true,
    pdftitle           = {test-005: spmoney, tabla de monedas},
  }
\begin{document}

\section*{A. Símbolos y alias}

% doc: «peso / pesos»
1. Peso, símbolo: $\spmoney{500 \$}$.

% doc: «peso / pesos»
2. Peso, alias peso: $\spmoney{500 peso}$.

% doc: «peso / pesos»
3. Peso, alias pesos: $\spmoney{500 pesos}$.

% doc: «euro / euros»
4. Euro, símbolo: $\spmoney{500 €}$.

% doc: «euro / euros»
5. Euro, alias euro: $\spmoney{500 euro}$.

% doc: «euro / euros»
6. Euro, alias euros: $\spmoney{500 euros}$.

% doc: «libra / libras»
7. Libra, símbolo: $\spmoney{500 £}$.

% doc: «libra / libras»
8. Libra, alias libra: $\spmoney{500 libra}$.

% doc: «libra / libras»
9. Libra, alias libras: $\spmoney{500 libras}$.

% doc: «yen / yenes»
10. Yen, símbolo: $\spmoney{500 ¥}$.

% doc: «yen / yenes»
11. Yen, alias yen: $\spmoney{500 yen}$.

% doc: «yen / yenes»
12. Yen, alias yenes: $\spmoney{500 yenes}$.

% doc: «yuan / yuanes»
13. Yuan, alias yuan: $\spmoney{500 yuan}$.

% doc: «yuan / yuanes»
14. Yuan, alias yuanes: $\spmoney{500 yuanes}$.

% doc: «won / wons»
15. Won, símbolo, con money-math=false: $\spmoney[money-math=false]{500 ₩}$.

% doc: «won / wons»
16. Won, alias won, con money-math=false: $\spmoney[money-math=false]{500 won}$.

% doc: «won / wons»
17. Won, alias wons, con money-math=false: $\spmoney[money-math=false]{500 wons}$.

% doc: «séquel / séqueles»
18. Séquel, símbolo, con money-math=false: $\spmoney[money-math=false]{500 ₪}$.

% doc: «séquel / séqueles»
19. Séquel, alias shekel, con money-math=false: $\spmoney[money-math=false]{500 shekel}$.

% doc: «séquel / séqueles»
20. Séquel, alias shekels, con money-math=false: $\spmoney[money-math=false]{500 shekels}$.

% doc: «séquel / séqueles»
21. Séquel, alias sequel, con money-math=false: $\spmoney[money-math=false]{500 sequel}$.

% doc: «séquel / séqueles»
22. Séquel, alias sequeles, con money-math=false: $\spmoney[money-math=false]{500 sequeles}$.

% doc: «rupia / rupias»
23. Rupia, símbolo, con money-math=false: $\spmoney[money-math=false]{500 ₹}$.

% doc: «rupia / rupias»
24. Rupia, alias rupia, con money-math=false: $\spmoney[money-math=false]{500 rupia}$.

% doc: «rupia / rupias»
25. Rupia, alias rupias, con money-math=false: $\spmoney[money-math=false]{500 rupias}$.

% doc: «rublo / rublos»
26. Rublo, símbolo, con money-math=false: $\spmoney[money-math=false]{500 ₽}$.

% doc: «rublo / rublos»
27. Rublo, alias rublo, con money-math=false: $\spmoney[money-math=false]{500 rublo}$.

% doc: «rublo / rublos»
28. Rublo, alias rublos, con money-math=false: $\spmoney[money-math=false]{500 rublos}$.

% doc: «lira / liras»
29. Lira, símbolo, con money-math=false: $\spmoney[money-math=false]{500 ₺}$.

% doc: «lira / liras»
30. Lira, alias lira, con money-math=false: $\spmoney[money-math=false]{500 lira}$.

% doc: «lira / liras»
31. Lira, alias liras, con money-math=false: $\spmoney[money-math=false]{500 liras}$.

% doc: «grivna / grivnas»
32. Grivna, símbolo, con money-math=false: $\spmoney[money-math=false]{500 ₴}$.

% doc: «grivna / grivnas»
33. Grivna, alias grivna, con money-math=false: $\spmoney[money-math=false]{500 grivna}$.

% doc: «grivna / grivnas»
34. Grivna, alias grivnas, con money-math=false: $\spmoney[money-math=false]{500 grivnas}$.

% doc: «centavo / centavos»
35. Centavo, símbolo: $\spmoney{500 ¢}$.

% doc: «centavo / centavos»
36. Centavo, alias centavo: $\spmoney{500 centavo}$.

% doc: «centavo / centavos»
37. Centavo, alias centavos: $\spmoney{500 centavos}$.

\section*{B. Códigos ISO en mayúsculas}

% doc: «peso chileno / pesos chilenos»
38. Peso chileno, código en mayúsculas: $\spmoney{500 CLP}$.

% doc: «peso mexicano / pesos mexicanos»
39. Peso mexicano, código en mayúsculas: $\spmoney{500 MXN}$.

% doc: «dólar estadounidense / dólares estadounidenses»
40. Dólar, código en mayúsculas: $\spmoney{500 USD}$.

% doc: «euro / euros»
41. Euro, código en mayúsculas: $\spmoney{500 EUR}$.

% doc: «libra / libras»
42. Libra, código en mayúsculas: $\spmoney{500 GBP}$.

% doc: «yen / yenes»
43. Yen, código en mayúsculas: $\spmoney{500 JPY}$.

% doc: «yuan / yuanes»
44. Yuan, código en mayúsculas: $\spmoney{500 CNY}$.

% doc: «won / wons»
45. Won, código en mayúsculas: $\spmoney{500 KRW}$.

% doc: «séquel / séqueles»
46. Séquel, código en mayúsculas: $\spmoney{500 ILS}$.

% doc: «rupia / rupias»
47. Rupia, código en mayúsculas: $\spmoney{500 INR}$.

% doc: «rublo / rublos»
48. Rublo, código en mayúsculas: $\spmoney{500 RUB}$.

% doc: «lira / liras»
49. Lira, código en mayúsculas: $\spmoney{500 TRY}$.

% doc: «grivna / grivnas»
50. Grivna, código en mayúsculas: $\spmoney{500 UAH}$.

\section*{C. Códigos ISO en minúsculas}

% doc: «peso chileno / pesos chilenos»
51. Peso chileno, código en minúsculas: $\spmoney{500 clp}$.

% doc: «peso mexicano / pesos mexicanos»
52. Peso mexicano, código en minúsculas: $\spmoney{500 mxn}$.

% doc: «dólar estadounidense / dólares estadounidenses»
53. Dólar, código en minúsculas: $\spmoney{500 usd}$.

% doc: «dólar estadounidense / dólares estadounidenses»
54. Dólar, alias dolar: $\spmoney{500 dolar}$.

% doc: «dólar estadounidense / dólares estadounidenses»
55. Dólar, alias dolares: $\spmoney{500 dolares}$.

% doc: «euro / euros»
56. Euro, código en minúsculas: $\spmoney{500 eur}$.

% doc: «libra / libras»
57. Libra, código en minúsculas: $\spmoney{500 gbp}$.

% doc: «yen / yenes»
58. Yen, código en minúsculas: $\spmoney{500 jpy}$.

% doc: «yuan / yuanes»
59. Yuan, código en minúsculas: $\spmoney{500 cny}$.

% doc: «won / wons»
60. Won, código en minúsculas, con money-math=false: $\spmoney[money-math=false]{500 krw}$.

% doc: «séquel / séqueles»
61. Séquel, código en minúsculas, con money-math=false: $\spmoney[money-math=false]{500 ils}$.

% doc: «rupia / rupias»
62. Rupia, código en minúsculas, con money-math=false: $\spmoney[money-math=false]{500 inr}$.

% doc: «rublo / rublos»
63. Rublo, código en minúsculas, con money-math=false: $\spmoney[money-math=false]{500 rub}$.

% doc: «lira / liras»
64. Lira, código en minúsculas, con money-math=false: $\spmoney[money-math=false]{500 try}$.

% doc: «grivna / grivnas»
65. Grivna, código en minúsculas, con money-math=false: $\spmoney[money-math=false]{500 uah}$.

\end{document}
